% AI-written, AI cross-reviewed; no human sign-off; not a claim that RH is proved
% RH Notes blueprint source excerpts for Z4.1
% Verbatim excerpts in the notation of the parent blueprint. The labelled editions and source ranges appear below.
% See the unit YAML for the current proof status. Earlier-round passages retain their original wording.
\section*{Z4.1. Normalized Hecke action on jets}

% Parent source SHA-256: 1af64922b98ed2310744411b1f18b25cc426c97220646321471836e1a3a409d6
% blueprint/ch_P.tex lines 4055 through 4790; sec:P-r4-z41, sec:P-r4-z41-kernel
\subsection{Integral jet estimates and normalized action (Z4.1)}\label{sec:P-r4-z41}
Let $v(x)=-\log_p\|x\|$, with $v(0)=+\infty$. Assume
\path{EntryHyp}: $\nu>0$, $p>2$ is prime, the field is ultrametric with
$\|p\|=p^{-1}$, all $u_s,\xi_{s,a},T_{c,i,a}$ are integral,
$e_{s,b}=0$ for $s\notin TT$, and $F$ belongs to the coefficient class
of $h=hhS(p,u)$. Thus there are integral $U_i$ such that, in degree
$d$, $F$ has the coefficients of $\sum_{i=0}^{d}p^{-i}h^iU_i$.
Put $a_s=\sum_b e_{s,b}$, $L_s=\sum_{c:s_c=s}l_c$,
$C_s=\#\{c:s_c=s\}$, $d=\sum_s(a_s+L_s)$, $\rho=1/(p^2-1)$ and
$\gamma=p\rho$. We prove the three displayed inequalities above
for the factorial pairing $z$ with the frame entry. The variables
at each of the $2\nu$ cyclic positions form an $r_0$ by $r_1$ matrix.
See \path{InClass} (\path{OddLogRelation.lean:76}) and
\path{lowOp}, \path{entryOp}, \path{EntryHyp}, \path{EntryStar}
(\path{Formalism/ThetaLogCoordinates/UniformThetaSubstitution/StatementE.lean:147,151,170,176}).

\emph{Terms visible to an entry of degree $d$.}
Write $Y_{s,a,b}$ for low and $X_{s,b}$ for top variables. Replace the
original variable by $Y_{s,a,b}+\mathbf1_{s\in TT}\xi_{s,a}X_{s,b}$.
The chain rule for the factorial pairing moves the entry's top linear
forms into this substitution, leaving $X^e Q(Y)$ on the operator
side, where $Q=\prod_c\det((T_cY_{s_c})_{*,I_c})$.
Take $E$ with $d<p^{2E}$ and truncate the two alternating
Frobenius series in $h$ there. All discarded terms have degree $>d$.
In each surviving power of exponent $p^t$, the mixed binomial
coefficients are divisible by $p$. Expanding the two factors in each
product of $h$ leaves the following four types, with integral coefficients:
\[
 A_\phi,\qquad B_k,\qquad Z,\qquad pY_0 .
\]
In $A_\phi$ the same column occurs at $s-2f$ and $s+1-2g$,
with exponents $p^{2f+1}$ and $p^{2g}$, respectively. Both positions
lie in $TT$ when its coefficient is nonzero. A term $B_k$,
$k=(s,f,b)$, uses $p^{2f+1}$ top letters at $(s,b)$ and one low
linear form in column $b$ at $s+2f+1$, with vector
$(\xi_{s,a}^{p^{2f+1}})_a$. The terms grouped in $Z$ contain at least
$p$ low letters: the only smaller positive exponent of a low factor is
the linear term from $g=0$, which has already contributed to $B_k$.
All binomial errors go into $pY_0$. This is coefficientwise through
degree $d$; the integral $U_i$ stay integral under the linear
substitution. See \path{aeval_entry} (\path{NormLeOfDvd.lean:116}),
\path{dAof}, \path{bBof} (\path{Formalism/ThetaLogCoordinates/UniformThetaSubstitution/StatementE.lean:106,118}),
\path{sort_prod}, \path{A_eq}, \path{B_eq}, \path{decomposition_of_hh}
(\path{EDecomp.lean:288,352,414,498}).

\emph{Factorials.} Consider a nonzero term in the multinomial expansion
of $p^{-i}h^iU_i$. Let $P_\phi,B_k,i_2$ be its multiplicities of
$A_\phi,B_k,Z$, write $P=\sum_\phi P_\phi$, $B=\sum_k B_k$, and set
$i_{s,f}=\sum_bB_{s,f,b}$. Each $pY_0$ pays for its own denominator.
If $b_k$ is the direction of the low linear form, expand
$Q(Y+\sum_k t_k b_k)$ and take the coefficient of $\prod_k t_k^{B_k}$.
Multiplication of the series by one low linear form is adjoint for the
factorial pairing to directional differentiation of the polynomial.
Iterating gives precisely this coefficient times $\prod_k B_k!$, by
the multinomial theorem applied to each monomial. The multinomial
coefficient choosing the $B$ factors includes the factor
$B!/\prod_k B_k!$. Hence $B!$ survives. A residual low multi-index
$n$ contributes its factorial $n!$, while the top multi-index contributes
$e!$ (multi-index factorials are products). All other coefficient factors
are integral. For a nonzero contribution,
\[
 |n|+B=\sum_sL_s,\qquad |n_s|\le L_s,\qquad pi_2\le |n|,
 \qquad v(\text{term})\ \ge\ v_p(e!B!n!)-P-B-i_2 .
\]
Homogeneity gives the first two degree assertions and each factor of $Z$
gives $p$ low letters for the third. The pure and mixed top exponents
consumed are at most $e$ in every coordinate, as required for a nonzero
coefficient. The ultrametric inequality permits us to use the weakest
bound among these finitely many terms. See \path{e_one_identity_of},
\path{e_one_integral_of} (\path{FinsuppStepInduction.lean:189,446}),
\path{valid}, \path{bnd}, \path{term_bound}, \path{TermBound_of}
(\path{NormNatCastLeOfDvd.lean:386,392,397,495}).

Let $q_p(m)$ be the base-$p$ digit length, with $q_p(0)=0$. Summing
$\lfloor m/p^j\rfloor$ over $j\ge1$ gives
$v_p(m!)=(m-\mathrm{digitsum}_p(m))/(p-1)\ge m/(p-1)-q_p(m)$.
Use $v_p(B!)\ge\sum_{s,f}v_p(i_{s,f}!)$ by multinomial integrality.
Since $p^{2f+1}i_{s,f}\le a_s$, there are at most $q_p(a_s)$ nonzero
$i_{s,f}$ for a given $s$, each with at most $q_p(a_s)$ digits.
The top and residual low arrays have $r_1$ and $r_0 r_1$ entries per
position. Their losses and the mixed loss total at most
\[
 \Delta=\sum_s\bigl(r_1q_p(a_s)+r_0r_1q_p(L_s)\bigr)
                +\sum_{s\in TT}q_p(a_s)^2 .
\]
The three factorial degrees sum to $d$. Since $1/(p-1)=\gamma+\rho$
and $i_2\le\sum_sL_s/p\le\gamma\sum_sL_s$, this gives the common bound
\[
        v(z)\ \ge\ \rho d+\gamma\sum_{s\in TT}a_s-P-B-\Delta
        \tag{*}
\]
for the multiplicities of a worst contributing term. A zero entry
satisfies all three bounds directly. See \path{deLoss}
(\path{StatementAB.lean:100}), \path{factorialValuationLowerBound},
\path{sum_lgN_le}, \path{gainFactorialInequality}
(\path{ABArith.lean:35,55,103}) and \path{EntryStar_of}
(\path{NormLeOfDvd.lean:309}).

\emph{How the three charges differ.} For fixed $(s,f)$ the directions
in $B_{s,f,b}$ have a common vector and target $t=s+2f+1$. A determinant
at $t$ using two such directions has two proportional replaced columns,
so it vanishes. Thus $i_{s,f}\le C_t$. More generally, for hits using
columns in $J$, multilinearity bounds the hits on a nonzero determinant
by $|\operatorname{im}I_c\cap J|$. Repeated original column indices give
the zero determinant. Expand the product of the minors to sum the
individual bounds. See \path{minor_hits} (\path{FinsuppStepInduction.lean:746})
and \path{HitsCapPos_of}, \path{HitsColumnsPos_of}
(\path{MinorProdBlockDegree.lean:36,58}).

For an $A_\phi$, orient its pair from the larger exponent to the smaller.
Write $e_{\rm hi}=\max(2f+1,2g)$ and $e_{\rm lo}=\min(2f+1,2g)$.
If the charge starts at $s$ and lands at $s'$, then
$s+e_{\rm hi}=s'+e_{\rm lo}$ modulo $2\nu$. Thus
$\delta=(s'-s)_{\bmod 2\nu}$ is the least positive residue of
$e_{\rm hi}-e_{\rm lo}$; it is odd and at most $e_{\rm hi}$.
The charged top multiplicity is at least $p^\delta$. Let $N_{s,f}$
count charges starting at $s$ with $\delta=2f+1$. Then
\[
 N_{s,f}+i_{s,f}\le a_s/p^{2f+1},\qquad
 N_{s,f}=0\quad\hbox{if }2f+1\notin D(s),
\]
where $D(s)$ consists of the distances to opposite-parity tops.
Subtract these two counts from $a_s/p^{2f+1}$. On $D(s)$ the
remainder is nonnegative. Outside $D(s)$ it is at least
$[a_s/p^{2f+1}-C_{s+2f+1}]_+$. Summing over $s,f$, use
$\sum_{f\ge0}p^{-(2f+1)}=\gamma$. The left side is
$\gamma\sum_s a_s-P-B$. The right side minus $\Delta$ is $G_A$.
Substitution in (*) yields $v(z)\ge G_A+\rho d$.
See \path{charge_exists}, \path{entryExcessBound_of}
(\path{Formalism/ThetaLogCoordinates/UniformThetaSubstitution/FrobeniusLiftNorms/EntryAOneParity.lean:55,177}) and \path{excessGainBound}
(\path{ABArith.lean:193}).

Every paired term consumes at least one even top letter. An even-origin
mixed term does so too, hence $P+\sum_{s\ {\rm even},f}i_{s,f}
\le\sum_{s\ {\rm even}}a_s$. The remaining, odd-origin mixed counts are
bounded by $\min(a_s/p^{2f+1},C_{s+2f+1})$. Applying (*) gives
$v(z)\ge G_B-\sum_{s\ {\rm even}}a_s+\rho d$, by \path{gainG}
(\path{Formalism/GrowthSums/StatementA.lean:104}). For the weighted inequality, mark the
top coordinates with $s$ even and $b\notin J$, or $s$ odd and $b\in J$.
Every pair has opposite parity and a common column, so it consumes a
marked letter. Each even-origin mixed term either consumes a marked
letter ($b\notin J$), or hits a column of $J$ in an odd minor
($b\in J$). The marked top budget and the distinct-column budget of the
odd minors add to exactly $b_J$. Thus $P+\sum_{s\ {\rm even},f}i_{s,f}
\le b_J$. Use the same odd-origin estimate and add $v(p^{b_J})=b_J$
to (*) to obtain $v(p^{b_J}z)\ge G_B+\rho d$. See
\path{entryCappedBound_of}, \path{EntryB_of}
(\path{EntryATwoAndB.lean:220,268}) and \path{cappedGainBound}
(\path{ABArith.lean:240}).

\paragraph{From the three entry estimates to the three place bounds (Z4.1).}
\label{par:r5-z41-bridge}
Fix a place $v$, a step $j+1\le m$, and a column choice $I$. Write
$H_{sc}$ for the height of column $c$ at position $s$.
In the field of rational functions in the frame coordinates $Y_{sia}$
use the multiplicative Gauss norm. Constants retain their norm,
each $Y_{sia}$ is integral, and evaluating a frame polynomial at these
indeterminates has norm equal to its largest coefficient norm. Indeed choose, in each nonzero polynomial, the least multi-index of a
coefficient of largest norm in an additive monomial order. At the sum
of the two chosen indices their product strictly dominates every
other summand: another pair attaining the two maxima would have both
indices at least the chosen indices and their sum could agree only for
the chosen pair. This proves multiplicativity and permits extension by
quotients (\path{gN_mul}, \path{gAbs_algebraMap},
\path{Formalism/GaussJetBounds/GaussJetProofs/GField.lean:97--135,168--193}). The
coefficientwise isometric embedding preserves the integral witnesses and
all coefficient identities of the scaled series class. In the resulting
minor product separate precisely those $(s,c)$ with $s\in TT$ and
$H_{sc}=1$. Their $1$-by-$1$ determinants are the linear forms
$\sum_aY_{s0a}X_{sa b}$, with $b=I_{sc}(0)$. Group them by $(s,b)$
and write $e_{sb}$ for the number in each group. Enumerate the
complementary minors by $c'$ and retain their original row and column
indices. Their product with the grouped powers is the entry operator.
Thus $\sum_b e_{sb}$ counts the top singletons, the capacity counts the
complementary columns, and their total height is the total number of
boxes minus the number of top singletons. These are exactly
$\mathrm{alN}=\mathrm{alG}$, $\mathrm{capN}=\mathrm{capG}$ and
$\mathrm{lsN}=\mathrm{lsG}$. Evaluation commutes with the finite
differential pairing, so each coefficient of the frame jet has norm
bounded by this one entry pairing. See \path{generic_elements_exist},
\path{EGaussK.reduce}
(\path{Formalism/ThetaLogCoordinates/UniformThetaSubstitution/FrobeniusLiftNorms/EGaussK.lean:16--32,38--101})
and \path{EGauss.minorAt_one}, \path{EGauss.entry_exists},
\path{EGauss.struct}
(\path{Formalism/ThetaLogCoordinates/UniformThetaSubstitution/FrobeniusLiftNorms/EGauss.lean:55--92,111--206}).

For the weighted bound set $J=\{b:b<i\}$, $i\ge1$, and let
$t_{sc}=|\operatorname{im}I_{sc}\cap J|$. The system exponent is
\[
 E_i=\sum_{s\ {\rm even},c}\bigl(\min(H_{sc},i)-t_{sc}\bigr)
       +\sum_{s\ {\rm odd},c}t_{sc}.
\]
Here $t_{sc}\le\min(H_{sc},i)$. An even top singleton contributes
$1$ to $b_J$ exactly when its index is outside $J$, hence contributes
$1-t_{sc}$ to both sides. An odd top singleton contributes $t_{sc}$
on each side. Each remaining odd minor contributes $t_{sc}$ on
each side. A remaining even minor contributes zero to $b_J$ and a
nonnegative amount to $E_i$. This proves $b_J\le E_i$ column by
column. As $|p|=p^{-1}<1$, the estimate with $p^{b_J}$ implies the
one with $p^{E_i}$. Repeated column indices give a zero determinant.
Otherwise counting image indices is the same as counting their
preimages, so the scaled differential polynomial is exactly
$p^{E_i}$ times the unscaled one. Moreover the jet in system variables
is obtained from the local frame jet by an injective variable renaming.
A coefficient outside the image is zero and one inside is the
coefficient of its unique preimage.
See \path{EGauss.struct}, \path{GaussBK_of}
(\path{Formalism/ThetaLogCoordinates/UniformThetaSubstitution/FrobeniusLiftNorms/EGauss.lean:207--246},
\path{Formalism/ThetaLogCoordinates/UniformThetaSubstitution/FrobeniusLiftNorms/EGaussK.lean:166--184});
\path{ea_PIvS_eq}, \path{ea_brScale}, \path{eb_jet}, \path{eb_brJet}
(\path{Formalism/SharpRestriction/GrowthBounds/MultisetValuation/EaApplyDScale.lean:45--56,98--139},
\path{Formalism/SharpRestriction/GrowthBounds/MultisetValuation/EbBrJet.lean:24--48}).

Identify $\operatorname{Fin}(n'_v)$ with $\mathbf Z/(2\nu)$ by least
residue, $n'_v=2\nu$. At a top there are
$\nu_{s,1}-\nu_{s,2}$ singletons. At other positions put this count
equal to zero. Subtracting it from the first row length gives capacity;
subtracting it from the sum of the column heights gives the low size.
These formulas count columns and boxes even when an individual count
is zero. The residue identification preserves parity and sends
cyclic addition and differences to their counterparts, so it preserves
the excluded-distance sets and the even-top sum. The digit lengths
in the two loss formulas agree at positive integers and have the same
zero convention. Substitution in the definitions thus gives the three
exponents $g_A$, $g_B-\mathrm{alEven}$ and $g_B$ at the actual
system step $j+1$.
See \path{ea_alG}, \path{ea_capG}, \path{ea_lsG}, \path{ea_deLoss},
\path{gain_of_capacity_loss}, \path{ea_gainG}, \path{ea_sumEven},
\path{ea_brNum}
(\path{Formalism/SharpRestriction/GrowthBounds/MultisetValuation/GainOfCapacityLoss.lean:13--48,74--145,146--158,173--257}).

For the series class let $\epsilon$ insert this place into the full
coordinate set. Restriction to this place sends a series $F$ to
the coefficients $F_{\epsilon_*n}$. Injectivity of $\epsilon$ gives
a bijection on all decompositions of $\epsilon_*n$ as a sum of two
multi-indices supported there. Every such decomposition is supported
there, so restriction commutes with multiplication. It preserves
total degree, hence commutes with rescaling by $X\mapsto\varpi X$.
The monomial in $s\mathrm{SA}$ at position $s-2l$ and exponent
$p^{2l}$ restricts to the corresponding monomial of $s\mathrm{Pos}$
at exponent $(p^2)^l$. Summing and multiplying therefore sends
$hhA$ to $hhS$. Restrict the integral witnesses for each power in
the class identity; their coefficients remain in the unit ball.
This gives the scaled entry class used above.
See \path{ea_resV_mul}, \path{ea_resV_atPi}, \path{ea_atPi_subst},
\path{ea_res_sSA}, \path{ea_res_hhA}, \path{ea_brClass}
(\path{Formalism/SharpRestriction/GrowthBounds/MultisetValuation/EaDegreeMapDomain.lean:16--81,84--102,118--172}).

There is a normed extension with the required root and the correct
original integral ring $R=Y.O_c$. At a place $e_O>0$ and $p$ is
associated in the base valuation ring to a positive power of its
uniformizer. If $p^{-1}\in R$, intersection with the base field
would give an inverse of $p$ in the base valuation ring, contradicting
that association. For $x\ne0$ put
$v(x)=\sup\{a/b:b>0,\ x^bp^{-a}\in R\}$.
The rank-one condition gives a member and bounds the set above. Multiplying
two tests after clearing denominators gives $v(xy)\ge v(x)+v(y)$.
For every rational $a/b$, either its test for $x$ or the inverse
test for $x^{-1}$ is integral. A rational strictly between $v(x)$
and $-v(x^{-1})$ would violate that dichotomy; applying the product
bound to $xx^{-1}=1$ proves equality for inversion and then for
multiplication. For $x\notin R$ the rank-one condition gives
$(x^{-1})^n/p\in R$ with $n>0$, so $v(x)<0$, whereas $x\in R$
passes the test $0/1$. Factoring $x+y=y(1+x/y)$ when $x/y\in R$,
and interchanging $x,y$ otherwise, proves the ultrametric inequality.
Since $p^k\in R$ exactly for $k\ge0$, $v(p)=1$. The resulting
norm $p^{-v}$, with value zero at zero, has unit ball $R$. Complete
and use the extending spectral norm on the algebraic closure of the
completion. Its restriction to the original field is isometric.
The polynomial $T^{p^2-1}-p$ now has a root $\varpi$, and
$|\varpi|=p^{-1/(p^2-1)}$. See \path{nm_pinv_of_place}, \path{nm_ext},
\path{nm_restNorm_of_place}
(\path{Formalism/SharpRestriction/GrowthBounds/MultisetValuation/NmPinvOfPlace.lean:14--62},
\path{Formalism/SharpRestriction/GrowthBounds/MultisetValuation/NmRestNormOfPlace.lean:15--21});
\path{nm_v}, \path{nm_v_mul}, \path{nm_v_nonneg_iff}, \path{nm_v_p},
\path{nm_v_add}, \path{nm_extG}
(\path{Formalism/SharpRestriction/GrowthBounds/MultisetValuation/RationalLowerBounds.lean:30--84,92--106,145--173,181--238},
\path{Formalism/SharpRestriction/GrowthBounds/MultisetValuation/NmNEqZero.lean:64--85,97--125}).

Finally choose an arbitrary factor $U$ whose parameter expression has
integral coefficients. The scaled-series calculation above applies to
this $U$ after coefficient extension: the Witt coefficients and
$\varpi$ lie in the unit ball, hence so does the subring they generate.
It gives a unit-ball coefficient system $u$ and the scaled class just
restricted. Apply the appropriate jet bound, cancel the homogeneous
factor $\varpi^d$ as in the preceding Gauss descent, and denote the
coefficient obtained before field extension by $x$. Coefficient
extension commutes with the finite differential pairing, so
$|x|\le p^{-g}$ in the extension. For $b>0$ and $a\le bg$,
$|x^bp^{-a}|\le p^{-bg+a}\le1$. The unit-ball identity puts
$x^bp^{-a}$ in $Y.O_c$, precisely the test $Y.\mathrm{geC}(x,g)$.
The argument applies to every $U$, column choice and monomial. Taking
the two unscaled operators and the $i$-scaled one gives the stated
three place bounds. See \path{ea_geC_of_norm}, \path{ea_Rp_le_ball},
\path{ea_lowN}, \path{ea_A1U}, \path{ea_A2U}, \path{ea_BU}
(\path{Formalism/SharpRestriction/GrowthBounds/MultisetValuation/EaGeCOfNorm.lean:53--110,198--239});
\path{first_theorem_place_bound}, \path{second_lemma_place_bound} and
\path{hecke_place_bound}
(\path{Formalism/SharpRestriction/GrowthBounds/MultisetValuation/FirstTheoremPlaceBound.lean:20--30}).

\paragraph{Expansion at one place and multiplication over the places.}
Take the residue characteristic $p\ge5$, an even cycle of length
$n'\ge2$, the prescribed Frobenius-periodic trace-dual bases, a Witt
coefficient $c$ with $\sigma^2c=c$, $\varpi^{p^2-1}=p$, and a
zero-constant inverse $\lambda_F(\mu)=X$. Let $R$ be the subring
generated by the Witt coefficients and $\varpi$ in the extension
field (\path{Rp}, \path{FormalGroupSystem.lean:194}). Put $q=p^2$. In the scaled logarithm
$\lambda_F$, the coefficient at $q^j$ is
$\varpi^{q^j-1}/p^j=p^{1+q+\cdots+q^{j-1}-j}$, with exponent zero
at $j=0,1$. It is integral and divisible by $p$ for $j\ge2$.
Solving for the coefficient of $\mu$ in increasing degree gives an
integer recursively from lower coefficients since the linear term of
$\lambda_F$ is $1$. Modulo $p$ the inverse satisfies $\mu+\mu^q=X$.
Hence
\[
               \mu(X)\equiv\sum_{j\ge0}(-1)^jX^{q^j}\pmod p .
\]
Each coefficient in this identity is a finite calculation. See
\path{sx_lam_int}, \path{sx_mu_alt_int}, \path{sx_mu_alt}
(\path{SxMuAltInt.lean:13,187,332}).

Let $S_{t,a,b}$ denote the alternating sum in the variables at positions
$t-2j$ with exponents $p^{2j}$. On even positions the trace-dual basis
identity says that the two square basis matrices multiply to the
identity. Reversing them gives
$\sum_r\sigma^t(\omega_r^*)\sigma^{t'}(\omega_r)=\delta_{tt'}$
for $t,t'$ even. Applying $\sigma$ gives the odd--odd case.
In the residue ring the binomial theorem gives
$(\sum x)^p=\sum x^p$; Witt Frobenius reduces to the $p$th power.
Consequently the alternating inverse applied to each linear basis
coordinate is congruent modulo $pR$ to a sum of the $S_{t,a,b}$
with respective coefficients $\sigma^t(\omega_r^*)$ on the even
positions or $\sigma^t(\omega_r)$ on the odd ones. Indeed, in the
$j$th term shift $t$ by $2j$ and replace the $p^{2j}$th power of
the basis coefficient by its $2j$th Frobenius iterate. The bases
are periodic with period $n'$. In the product of one coordinate
and the $p$th power of the other, the latter Frobenius changes its
index by one, and the sum over $r$ vanishes unless the two indices
match after this shift. The two ordered products therefore give
\[
 \sum_{s,a,b}u_s S_{s,a,b}^{p} S_{s+1,a,b}
       =hhS(p,u),\qquad
  u_s=\begin{cases}c&s\ {\rm odd},\\-\sigma c&s\ {\rm even}.
       \end{cases}
\]
The equal sign here identifies the surviving cyclic polynomial in
the series coordinates; the substituted original difference equals
this series plus $pY$ for integral $Y$. See \path{sx_dual_even},
\path{sx_dual}, \path{frobenius_congruence_witt}, \path{sx_frob_sum}
(\path{FrobeniusCongruenceWitt.lean:13,26,80,102}) and
\path{restTwoSums_of} (\path{TripleSumComm.lean:205,286,298}).

In $G_{c,\sigma c}$ remove the leading expression
$(c X Y^p-\sigma c\,Y X^p)/p$ before scaling. In each orientation
the remaining terms have degree $q^i+p q^j$, denominator
$p^{i+j+1}$, $0\le i\le j$, $(i,j)\ne(0,0)$. After scaling both
variables by $\varpi$, their powers of $\varpi$ are at least $p+1$:
at $(0,1)$ the excess over $p+1$ is $(p-2)(q-1)$, increasing $i$
by $1$ changes it by $(q-1)(q^i-1)$ and increasing $j$ by $1$
changes it by $(q-1)(p q^j-1)$. The remaining series thus has the
form $\varpi^{p+1}W$ with $W$ integral and of zero constant.
After substituting the inverse coordinates, the negative exponent is
$-(\varpi^{p+1}/p)hhS+\varpi^{p+1}R_1$, with $R_1$ integral of zero
constant. Its second exponential is integral, since
$v_p(m!)\le m/(p-1)$ and $\varpi^{p+1}$ has valuation $1/(p-1)$
(the integer prime-to-$p$ factors of $m!$ are Witt units).
Renaming variables embeds each place in the disjoint union of places.
The sum of the local exponents has exponential equal to the product
of their exponentials, coefficient by coefficient, and the product of
their integral factors is integral. See
\path{exists_atPi_G1_eq_C_pow_mul_integral_series}
(\path{ExistsAtPiG1EqCPowMulIntegralSeries.lean:61}),
\path{gcCondition_of_twoSums} (\path{TwoSumsCondition.lean:174}),
\path{restGcA_of} (\path{ExpansionAllPlacesOfPlace.lean:173})
and \path{restSExpA_of} (\path{DworkExponentialRestrictions.lean:118}).

For a general $\beta$ satisfying \path{BeProp},
Section~\ref{sec:P-r4-z41-kernel} proves that
$\exp(\operatorname{inParam}(\beta)-G_{c,\sigma c})$ belongs to
$1+(X,Y)W[[X,Y]]$, over the original Witt coefficient ring $W$,
with $\sigma^2c=c$. A product of these units over the input
coordinate pairs and the unit from the given equivalence of $\Phi$
supply the correction for the exponential of $-\Phi$. Their inverses
also have integral coefficients by the geometric series in a
zero-constant term. Substitution in the integral zero-constant
coordinates preserves integrality. For the admitted kernel's
additional factor $U$, its expression in parameters is integral over
$R_1$. For each of the linear forms $L_i$ of the chosen Witt bases,
the integral coordinates $M_i=\varpi\mu(L_i)$ satisfy
$\log_F(M_i)=\varpi L_i$ by the inverse identity. Consequently the
scaled substitution on $U$ is the substitution of $M_i$ into its
integral parameter expression. It stays integral in every allowed
target subring $R'$ containing $R$ and the image of $R_1$. This is
\path{restBeta} (\path{BiadditiveSecondVariable.lean:137}),
\path{E0} (\path{IntegralExpRename.lean:61}),
\path{restSExp_of}, \path{restSExpA_of}
(\path{DworkExponentialRestrictions.lean:95,118}),
\path{genInt}, \path{restSInt} (\path{InParamMap.lean:58,124})
and \path{restSIntA} (\path{AtPiEqRescale.lean:64}).

Finally set $c_0=\varpi^{p+1}$, and let $h'=hhS+pH$ be any allowed
lift in $R'$ with zero constant. Then $H$ has zero constant as well.
The identity
\[
   \exp(-c_0 hhS/p)=\exp(-c_0h'/p)\exp(c_0H)
\]
absorbs an integral factor, since $c_0^i/i!$ belongs to $R$.
Writing the resulting integral prefactor as $V'$, take
$U_i=(-c_0)^i V'/i!$ in the definition of \path{InClass}.
Only $i\le d$ can matter in degree $d$. At multiple places the
corresponding factor is $V'\prod_v(-c_0)^{i_v}/i_v!$,
with $0\le i_v\le d$ for every $v$. This proves the required
coefficients for each allowed lift and after extension of subrings.
See \path{exp_change}, \path{SConcl_of}
(\path{KernelExpansionOfRestrictions.lean:94,116}),
\path{inClass_of_exp}, \path{inProd_of_exp}
(\path{InClassOfExp.lean:143,163}) and \path{all_places_of_restrictions}
(\path{AllPlacesOfRestrictions.lean:37}).

\paragraph{Why the normalized Hecke operator preserves the lattice.}
At an elementary edge over the principal ideal ring, take the Smith
normal form for the inclusion of source in target. Each elementary
divisor divides the irreducible $p$, so is a unit or is associated to
$p$. If the residue field has $p^{n'}$ elements, the quotient has
order $(p^{n'})^i$; the product of the cyclic factor orders therefore
gives exactly $i$ nonunit divisors. Permute these first. The source
basis matrix is $g$ and the target is the column span of
$gD$, $D=\operatorname{diag}(p^{-1}I_i,I)$.
See \path{lc_snf} and \path{lc_main}
(\path{Formalism/SharpRestriction/GrowthBounds/LatticeSystems.lean:2036--2042,2098--2125,2149--2152,2190--2214}). Integral changes
of basis on a fixed lattice lie in $\mathrm{GL}(\mathcal O)$:
express each basis in the other to obtain inverse integral matrices.
Their row actions preserve the integral Weyl lattice and their
determinants are units, so the twisted lattice is independent of
this choice. See \path{lc_main} (\path{LatticeSystems.lean:2098}),
\path{dl_indep} and \path{dl_pos} (\path{GShiftRowGenMonomial.lean:376,443}).

For an antitone polynomial weight $\eta$ put $N_i=\sum_{t<i}\eta_t$.
A divided row operator picking $a$ letters moves precisely $a$
row indices in each surviving monomial, so applying a fixed word to
the highest-weight vector gives a homogeneous polynomial for the
total degree $k$ in the first $i$ rows. Every supported monomial
satisfies $k\le N_i$: in a product of column minors of shape $\eta$
a column of height $h$ can contribute at most $\min(h,i)$ such
letters. Sum over the columns to count the boxes of the first $i$
rows. This continues to hold under divided row operations: expand
the minors after adding a parameter times one row to another and
extract the desired parameter power. Every term is still a product
of minors with the same column heights. The action of $D$ on this
word multiplies by $p^{-k}$; multiplication by $p^{N_i}$ leaves the
integral multiple $p^{N_i-k}$ of the same generator. The words
span, so the inclusion holds on the lattice. See \path{g_applyWord},
\path{dom_dualWeyl}, \path{diag_monomial}, \path{weyl_diag}
(\path{GShiftRowGenMonomial.lean:165,176,237,264}).

For an integral antitone weight put $\eta_t=\lambda_t-\lambda_{\rm last}$.
The determinant twist contributes $p^{i\lambda_{\rm last}}$
for $i$ in range. This changes $N_i$ to $\sum_{t<i}\lambda_t$.
A non-antitone weight gives the zero lattice in this construction.
The embeddings and the Frobenius powers carry the diagonal entries
$p^{-1},1$ to themselves. Tensor the single-factor inclusions over
the cycle, adding their exponents to obtain $A_i$, and take the
identity at the other places. Summing over the finite neighbors
preserves the target submodule. For $1\le i$ in the admissible
range and $j+1\le m$, the even weights give exactly the even part
of the normalizing exponent
\[
 N'(j+1,i)=\sum_{s\ {\rm even}}\sum_{a=1}^{i}\lambda_I(s,a)
              +i\sum_{s\ {\rm odd}}w(j+1,s)
\]
(\path{lamI}, \path{Nprime},
\path{Formalism/SharpRestriction/GrowthBounds/FormalGroupSystem.lean:23,72--74}).
Each of the at most $i$ odd
weights at a position $s$ is bounded by $w(j+1,s)\ge0$, giving
the odd part's upper bound. Thus $N'-A_i\ge0$ and multiplying the
already integral sum by $p^{N'-A_i}$ yields the prescribed normalized
operator. See \path{dl_P} (\path{GShiftRowGenMonomial.lean:327}),
\path{Tminus_mem}, \path{Ar_lift_le}, \path{TmInt},
\path{restricted_hecke_integral_of_edges}, \path{restWeylEdge}
(\path{RestrictedHeckeIntegralOfEdges.lean:35,74,136,157,234}).

\paragraph{The product entry and the single denominator loss.}
The admitted total kernel is $\prod_v F_v$ times a factor $U$
integral in parameters, by the exponential splitting above. A general
such $U$ can involve several places. Expand its parameter expression
in monomials. At fixed jet degree only finitely many contribute, since
every substituted coordinate has zero constant. In a given monomial
collect the coordinates over each $v$ as $U_v$, expressed in the
original variables by the inverse parameter series. Its parameter
expression is precisely the corresponding ordinary monomial, hence is
integral. The local estimate applies to $F_vU_v$ by its quantifier over
all such factors. The variables and derivative operators now split by
place: a coefficient of the product pairs with a product of local
derivatives, with the multi-index factorial splitting into the local
factorials. Consequently the entries have valuation at least
$\sum_v g_v$ for any local rational bounds $g_v$. Summing back the
monomials with integral coefficients preserves the same bound. This
proves the product estimate with the possibly mixed prefactor.
See \path{en_places_add}, \path{restPlaces_of}
(\path{KernelProductForm.lean:51,118}) and \path{restEnt_of}
(\path{RestEntOf.lean:62}).

By the admitted class formula, each component of the operator after
scalar extension is a finite sum over $x$ of
$c_K w_x^{-1}\sum_y K_{x,y}F_{x,y}$.
At an interior level $w_x$ is the stabilizer order for its chosen
representative; the last level has the two finite index sets and positive
weights of the last class formula. In both cases $c_K$ is integral
and $v_p(w_x)\le d_j$. The admitted frame condition at the chosen
related lattices converts the product entry bound into membership
of each normalized kernel value in the lattice of level $g$.
Thus $p^{d_j}c_K/w_x$ is integral for each $x$, and the two finite
sums cost only $d_j$ altogether. To descend, express the extended vector as a finite sum
$\sum a_k z_k$ of lattice vectors with coefficient lower bound
$g-d_j$. If the $z_k$ are dependent over the original field, choose a
relation and a nonzero coefficient of least valuation in it. Divide
the relation by that coefficient and eliminate its vector. Every
ratio is integral, hence the remaining coefficients still have the
same lower bound. Repeat until the vectors are independent. A
nonsingular coordinate minor of the remaining vectors then recovers
their coefficients from the original vector over the original field.
The bound on the recovered coefficients gives precisely membership
in the original scaled lattice (\path{descent_core},
\path{DescentCore.lean:98}; \path{descent}, \path{InLatCSmulC.lean:54}). For a list of distinct selected places, expand the Hecke
neighbor sums successively. At an unselected place the index is zero
and its scalar equals one. Each term in the expansion therefore
satisfies precisely the simultaneous edge relation, and the same bound
applies to the sum. Equivariance gives the asserted global lattice
condition. Passing from rational lower bounds to real endpoints is the
ceiling argument given above. See \path{KSum}, \path{ksum_entry},
\path{kernel_sum_of_class_data}, \path{entry_mem}, \path{TtL_mem},
\path{Eqv_TtL}, \path{formB_of}
(\path{KernelSumOfClassData.lean:178,185,197,206,92,244,255}) and
\path{formB_lim}, \path{form_of_place}
(\path{RamifiedLatticeOfChainEstimates.lean:65,90}).

\emph{Status: complete for Z4.1 under the stated admission hypotheses.}
Section~\ref{sec:P-r4-z41-kernel} supplies the biadditive normal form
used here. \textit{[Hypothesis, discharged at
\path{admissible_at_each_site} in \path{site_cmpLite_siteAdm}.]}
This application to the selected site uses the thirteen conjuncts of
\path{admitted}; the site realization belongs to Z4.2
(\path{Formalism/SharpRestriction/GrowthBounds/FormalGroupSystem.lean:771--773},
\path{Development/AbelianVarietyTheory/Proofs/SatakeParameters/JointSiteChainComparedCarriesAdmittedSystem.lean:29--46,77--89}).
The column proposition~\ref{prop:P4-columns} also requires the two
prefactor inequalities, the neighbour count and the inverse convolution
(\path{vGeK_cl}, \path{Cii}, \path{CountLoc}, \path{heckeCondition_of_adm};
\path{Formalism/SharpRestriction/GrowthBounds/MultisetValuation/TheoremOfAlternateForm.lean:524,538},
\path{Formalism/SharpRestriction/GrowthBounds/MultisetValuation/CountLoc.lean:9},
\path{Formalism/SharpRestriction/GrowthBounds/MultisetValuation/HeckeConditionOfAdm.lean:75}).


\paragraph{Classification of the biadditive local kernel.}\label{sec:P-r4-z41-kernel}
We prove the normal-form input \path{RestBeta} of the preceding
expansion. Work over $K=\operatorname{Frac}W(\overline{\mathbf F}_p)$,
$p\ge5$, with integral ring $O$, $v(p)=1$ and Witt Frobenius $\sigma$.
Write $q=p^2$, $N=q-1$ and
\[
 L(T)=\sum_{i\ge0}T^{q^i}/p^i,\qquad
 U(T)=L(T^p)/p,\qquad
 A(X,Y)=\sum_{0\le i\le j}
      X^{q^i}Y^{p q^j}/p^{i+j+1}.
\]
Here $U(T^p)=L(T)-T$. Write $B\sim B'$ when $B-B'$ has zero constant
and $\exp(B-B')$ is integral. For the input $b$ of \path{BeProp}, put
$B(X,Y)=b(L(X),L(Y))$. We shall obtain $c\in O$, $\sigma^2c=c$,
with $B\sim cA(X,Y)-\sigma(c)A(Y,X)$
(\path{RestBeta}, \path{DworkExponentialRestrictions.lean:15};
\path{BeProp}, \path{FormalGroupSystem.lean:125--157};
\path{lgF}, \path{g0F}, \path{Aser}, \path{Gser},
\path{OddLogRelation.lean:13--30}).

\emph{Two coefficient tests.} If $\partial_X h$ is integral, replacing
an integral zero-constant argument $x$ by $x+p r$ with $r$ integral of
zero constant changes $h$ by an element of $pO[[\mathbf X]]$, for
integral zero-constant arguments in the other variables as well.
For a term $h_n X^n$ the contribution using $r^k$, $k\ge1$, carries
$(nh_n)\binom{n-1}{k-1}p^k/k$, which belongs to $pO$. This proves
the assertion coefficientwise. In particular $L'$ is integral with
constant $1$. The unique $F=L^{-1}(L(X)+L(Y))$ is integral. Indeed it
has rational coefficients. If its truncation $F_{<d}$ is $p$-integral,
then $F_{<d}^q\equiv F_{<d}(X^q,Y^q)\pmod p$ by reduction to
$\mathbf F_p$. The test for $L$ gives
$L(F_{<d}^q)\equiv L(F_{<d}(X^q,Y^q))\pmod p$.
In degree $d>1$ the latter agrees with $L(X^q)+L(Y^q)$, since the
possible contributing degree of the original discrepancy is $d/q<d$.
Using $L(T)=T+L(T^q)/p$, the degree-$d$ discrepancy of $F_{<d}$
is integral. Its negative is the necessary degree-$d$ correction;
the induction starts at $X+Y$. Uniqueness gives $\sigma F=F$ and
$F(\zeta X,\zeta Y)=\zeta F(X,Y)$ for $\zeta^N=1$, as
$L(\zeta X)=\zeta L(X)$ (\path{divided_derivative},
\path{IntegralPderivLgF.lean:101,176}; \path{formal_law_exists},
\path{formal_law_symmetric}, \path{CoeffSubstLgFOfDegreeOne.lean:306,312}).

Set $\Phi h=(\sigma h)(X_1^p,\ldots,X_s^p)$ and
$Dh=\Phi h/p-h$. For $h(0)=0$,
\[
              \exp h\ \hbox{integral}\quad\Longleftrightarrow\quad
              Dh\ \hbox{integral}.
\]
Indeed, for $E=\exp h$, $\Phi E/E^p=\exp(pDh)$. If $E$ is integral
then $\Phi E\equiv E^p\pmod p$, and taking the logarithm of
$1+pT$ gives coefficients in $pO$. Conversely $Dh$ integral gives
$\Phi E/E^p\in1+pO[[\mathbf X]]$ by the exponential series.
Induct on degree in $E$. Dividing out its integral truncation of
degrees $<d$ preserves this ratio condition. If the remaining leading
homogeneous part is $C_d$, the degree-$d$ part of the ratio is $-pC_d$,
so $C_d$ is integral. This also proves closure of $h\sim0$ under
division by $N$, since $N$ is a $p$-adic unit fixed by $\sigma$
(\path{dwork_multivariable}, \path{Formalism/ThetaLogCoordinates/UniformThetaSubstitution/FrobeniusLiftNorms/LocalComponentGL.lean:239,448};
\path{LamModule_of}, \path{ScalingSwapLemmas.lean:314}).

\emph{Support and the linear coefficient.} Write
$\delta f=f(F(X,Y))-f(X)-f(Y)$, and $\delta_X G$ for this operation
on the first variable of a two-variable series. The two additive
clauses for $b$ give $\delta_X B\sim0$, $\delta_Y B\sim0$ after
substituting $L$ in each variable. The conjugacy clause gives
$B(\zeta X,Y)\sim B(X,\zeta^pY)$: each $N$th root is the
Teichm\"uller lift of a root in $\mathbf F_{p^2}^\times$, with
Frobenius its $p$th power. The average
\[
 G=N^{-1}\sum_{\zeta^N=1} B(\zeta X,\zeta^{-p}Y)
\]
is equivalent to $B$, and both additive defects still have integral
exponentials by the symmetry of $F$. Summing powers of the primitive
root shows that $G_{ij}=0$ unless $i,j\ge1$ and $i\equiv pj\pmod N$.
Swapping variables preserves this condition since $p^2\equiv1\pmod N$
(\path{sx_root}, \path{sx_betaHyp}, \path{BiadditiveSecondVariable.lean:61,130};
\path{StepZero_of}, \path{ScalingSwapLemmas.lean:340}).

Here is the one-variable calculation needed twice below. If $f(0)=0$
and $\delta f$ is integral, then for some $a\in O$, $b'\in K$,
\[
                    f=aU+b'L+R,\qquad R\in O[[T]].
\]
First $\delta U=(L(F^p)-L(F(X^p,Y^p)))/p$ is integral by the
substitution test and the congruence of the two arguments modulo $p$.
Subtracting an integral multiple of $U$ or any multiple of $L$
therefore preserves the assumption. Start with the remainder
$f-[T]f\cdot L$. If its coefficients below $n>1$ are integral, the
mixed coefficients of total degree $n$ in its coboundary imply
$\binom nk [T^n]R\in O$, $0<k<n$, since the linear part of $F$ is
$X+Y$. If $n$ is not a power of $p$, the base-$p$ digits of $n$ give
a proper $k=p^s$ with $\binom nk$ nonzero modulo $p$ (expand
$(1+T)^n$ in characteristic $p$). If $n=p^k$, $k\ge1$,
$\binom{p^k}{p^{k-1}}$ has valuation one, by
$\binom{p^k}{p^{k-1}}=p\binom{p^k-1}{p^{k-1}-1}$.
So the exceptional coefficient $r_n$ satisfies $p r_n\in O$.
At $n=p q^j$ subtract $p^{j+1}r_n U$, and at $n=q^{j+1}$
subtract $p^{j+1}r_n L$. These corrections kill $r_n$, preserve
integrality at all smaller degrees and change each scalar by an
element of $p^j O$. The two resulting sequences converge by
completeness and each limiting remainder coefficient belongs to the
closed ring $O$. See \path{DeIntegral_of} (\path{IntegralPderivLgF.lean:191}),
\path{correct}, \path{exists_limit} and \path{LemmaOne_of}
(\path{NotDvdChoosePowMul.lean:16--62,195--276,297--328,428--539}).

Put $f(Y)=[X]G(X,Y)$. Logarithmic differentiation of the integral
unit $\exp(\delta_Y G)$ in $X$ at $X=0$ gives $\delta f$ integral.
Differentiating $\exp(\delta_XG)$ in the added variable at zero gives
$\partial_X G/L'(X)-f(Y)$ integral, since differentiating the group
law gives $\partial_2F(X,0)=1/L'(X)$. In the preceding decomposition
of $f$, only exponents $j\equiv p\pmod N$ can occur in $f$, while
$L$ is supported on $j\equiv1$ and $U$ on $j\equiv p$.
Thus $f-cU$ is integral for the integral $c=a$. We have proved
$\partial_X G-cL'(X)U(Y)$ integral, and obtain $c'\in O$ likewise
from $G(Y,X)$ (\path{FirstDeriv_of}, \path{CobLin_of}, \path{StepI_of},
\path{ConstantCoeffLawFirstTwoOfThree.lean:265,337,411--473}).

\emph{The Frobenius correction.} Let $S=F(X,Y)$, $T=F(X^p,Y^p)$,
$Z$ the unchanged third variable and
$V=(\sigma G)(X,Z)-\sigma(c)L(X)U(Z)$. Its first derivative is
integral. Since $S^p\equiv T\pmod p$, the first coefficient test gives
$V(T,Z^p)-V(S^p,Z^p)\in pO[[X,Y,Z]]$. Direct substitution using
$L(S^p)-L(T)=p\delta U$ and $U(Z^p)=L(Z)-Z$ gives
\[
 \delta_X(DG)-\sigma(c)\delta U\, L(Z)
 =D(\delta_XG)-\frac{V(T,Z^p)-V(S^p,Z^p)}p-\sigma(c)\delta U\,Z.
\]
All three right-hand terms are integral, the first by the exponential
test. Applying the same identity with variables swapped shows that both
additive defects of
\[
 H=DG-\sigma(c)U(X)L(Y)-\sigma(c')L(X)U(Y)
\]
are integral. This is the commutation computation of \path{StepII_of}
(\path{PartialCoboundaryLinear.lean:225--245,320--416}).

The support of $H$ still has $i\equiv pj\pmod N$ and avoids the axes.
We give the slice argument that makes $H$ itself integral. Set
$u_i=[T^i]U$, $l_i=[T^i]L$. Both additive defects are now integral,
so the one-variable calculation on each row and column gives
\[
 H_{ij}\equiv a_j u_i+b_j l_i\equiv a'_i u_j+b'_i l_j\pmod O,
 \qquad a_j,a'_i\in O .
\]
Here $\pmod O$ means that the differences lie in the additive subgroup
$O$. Fix $i=p q^k$, and take $j=q^n$. The first expression has norm
at most $p^{k+1}$ independently of $n$; the second equals $b'_i/p^n$.
Hence $b'_{p q^k}=0$. Now fix $j$ and let $i=p q^k$ vary. The second
expression has norm at most $|u_j|$ and the first equals $a_j/p^{k+1}$,
so $a_j=0$. For $j$ not a power of $q$, take $i=q^k$. The second
expression stays bounded by $|u_j|$, giving $b_j=0$. For $j=q^n$,
$i=q^k$ is outside the support, since $1\not\equiv p\pmod N$.
Then $b_j/p^k\in O$ for every $k$, again yielding $b_j=0$.
This proves integrality of $H$ on the required support
(\path{LemmaTwo_of}, \path{LemmaTwoOf.lean:280--476};
\path{StepH_of}, \path{ConstantCoeffFirstTwoOfThree.lean:277--301}).

At bidegree $(1,pq^i)$ the derivative test says
$G_{1,pq^i}-c/p^{i+1}\in O$. At the same bidegree $\Phi G$ has
coefficient zero, so integrality of $H$ gives
$-G_{1,pq^i}-\sigma(c')/p^{i+1}\in O$. Thus $c+\sigma(c')$ is
in every $p^{i+1}O$ and is zero. Transposing gives
$c'+\sigma(c)=0$, hence $\sigma^2 c=c$. Put
$A_c=cA(X,Y)-\sigma(c)A(Y,X)$. In $DA_c$, shifting the indices by
applying $X,Y\mapsto X^p,Y^p$ and using $\sigma^2c=c$ gives
\[
                 DA_c=\sigma(c)U(X)L(Y)-cL(X)U(Y).
\]
For example at $(p q^i,q^j)$ the $-A_c$ term supplies
$\sigma(c)/p^{i+j+1}$ when $j\le i$, and the $\Phi A_c/p$ term
supplies it when $j>i$; the transposed coefficient is obtained with
factor $-c$. These cover the product supports. We have
$D(G-A_c)=H$ integral, so the exponential test gives $G\sim A_c$.
In the Witt field, membership $c\in O$ represents an actual Witt
vector and the embedding intertwines Frobenius. This proves the
normal-form input used in \path{restSExp_of} and \path{restSExpA_of};
the preceding one-place computation and its finite product then apply
to the original kernel
(\path{c_add_sigma_c_eq_zero}, \path{StepIV_of},
\path{ConstantCoeffFirstTwoOfThree.lean:338--419};
\path{BetaIsG_of}, \path{NormLiftLeOne.lean:32--60};
\path{restBeta}, \path{BiadditiveSecondVariable.lean:137--153};
\path{restSExp_of}, \path{restSExpA_of},
\path{DworkExponentialRestrictions.lean:95--135}).

